Yao et al.'s ReAct$\to$CoT-SC back-off trigger encodes a single failure
mode --- step exhaustion --- and, we show, misses a second, equally real
failure mode: premature confidence on thin evidence, which dominates
FEVER-style claim verification in our reproduction. Generalizing the
trigger is not free, however: the naive generalization (counting
retrieved evidence) does not work, contributing negligible or negative
precision lift over the base error rate on both domains we tested. An
entailment-verification signal, at the cost of one additional LLM call
per question, does work, and combined with the paper's original
step-exhaustion condition raises precision substantially on HotpotQA while
recovering some of FEVER's blind spot.

More broadly, this study is a case for validating a proposed signal
against an oracle label \emph{before} committing compute to a full
ablation: the negative result for evidence-count thinness was cheap to
obtain and, once obtained, changed which arms were worth running at full
scale. We report that negative result alongside the positive one, since a
generalization of a paper's mechanism that silently discards its own
failed hypotheses is less useful to future work than one that documents
what did not generalize and why.
